\documentclass[11pt]{article}
\usepackage[T1]{fontenc}
\usepackage{lmodern}
\usepackage[margin=1in]{geometry}
\usepackage{amsmath,amssymb,amsthm,mathtools,bm}
\usepackage{microtype,booktabs,array,enumitem}
\usepackage[hidelinks]{hyperref}
\usepackage{fancyvrb}
\hypersetup{pdfauthor={Sidney Holden},pdftitle={Exact Rank of Complex Hankel Tensors},pdfsubject={An exact-rank formula for TR-14}}
\setlength{\parskip}{0.35em}
\setlength{\parindent}{0pt}
\setlength{\emergencystretch}{2em}
\newtheorem{theorem}{Theorem}[section]
\newtheorem{lemma}[theorem]{Lemma}
\newtheorem{proposition}[theorem]{Proposition}
\newtheorem{corollary}[theorem]{Corollary}
\theoremstyle{definition}
\newtheorem{definition}[theorem]{Definition}
\newtheorem{example}[theorem]{Example}
\theoremstyle{remark}
\newtheorem{remark}[theorem]{Remark}
\newcommand{\C}{\mathbb C}
\newcommand{\Rk}{\operatorname{R}}
\newcommand{\SR}{\operatorname{R}_{\mathrm{sym}}}
\newcommand{\VR}{\operatorname{R}_{\mathrm V}}
\newcommand{\BR}{\underline{\operatorname{R}}}
\newcommand{\BSR}{\underline{\operatorname{R}}_{\mathrm{sym}}}
\newcommand{\BVR}{\underline{\operatorname{R}}_{\mathrm V}}
\newcommand{\rank}{\operatorname{rank}}
\newcommand{\Stab}{\operatorname{Stab}}
\newcommand{\Ann}{\operatorname{Ann}}
\newcommand{\spanC}{\operatorname{span}_{\C}}
\newcommand{\rad}{\operatorname{rad}}
\newcommand{\supp}{\operatorname{supp}}
\newcommand{\Spec}{\operatorname{Spec}}
\newcommand{\im}{\operatorname{im}}
\newcommand{\Id}{\operatorname{Id}}
\newcommand{\Pq}{\C[t]_{\le q}}
\newcommand{\ceil}[1]{\left\lceil #1\right\rceil}
\newcommand{\floor}[1]{\left\lfloor #1\right\rfloor}
\title{Exact Rank of Complex Hankel Tensors}
\author{Sidney Holden\\{\small Center for Computational Biology, Flatiron Institute}\\{\small Simons Foundation, New York, USA}}
\date{12 September 2026}
\begin{document}
\maketitle

\begin{abstract}
Let $H_{i_1\cdots i_m}=h_{i_1+\cdots+i_m}$ for $0\le i_j\le n-1$, over $\C$.
Write $D=m(n-1)$, and let $r$ be the rank of the middle Hankel catalecticant.
For $H\ne0$, choose a homogeneous binary apolar polynomial of smallest degree $r$,
and let $s$ be its number of distinct projective zeros.
We prove the exact formula
\[
\Rk(H)=\SR(H)=\min\{D-r+2,\ (m-1)r-(m-2)s\}.
\]
When the smallest-degree apolar polynomial is not unique, the first term is $r$,
so the formula is independent of this choice. This proves the universal exact-rank
equality requested in TR-14, including exceptional and mixed-multiplicity tensors.
The lower bound uses a contextual product-space inequality: an annihilating
Frobenius functional excludes nilpotents from the auxiliary stabilizer algebras,
leaving finitely many support partitions. The upper bounds are explicit symmetric
interpolation constructions. We also prove that all three border ranks equal $r$,
and that the maximum ordinary and symmetric Hankel ranks are
$(m-1)(n-1)+1$.
\end{abstract}

\section{Statement and conventions}

Throughout, the field is $\C$, $m\ge3$, $n\ge2$, $q=n-1$, and $D=mq$.
Tensor indices are zero-based. Thus the tensor considered here is exactly the
one in TR-14 after subtracting one from each index. Ordinary rank permits
arbitrary decomposable summands; symmetric rank permits only scalar multiples
of $v^{\otimes m}$. The question imposes neither a genericity assumption nor a
Vandermonde restriction.\footnote{Open Problems in Numerical Linear Algebra,
TR-14, ``Comon's exact-rank conjecture for Hankel tensors,'' accessed September 12,
2026. \url{https://github.com/ajt60gaibb/OpenProblemsInNLA/tree/main/tensor-computations/TR-14}.}

The problem originates in Nie--Ye, Conjecture 6.3.\footnote{J.~Nie and K.~Ye,
\emph{Hankel Tensor Decompositions and Ranks}, SIAM J. Matrix Anal. Appl.
40 (2019), 486--516. \url{https://doi.org/10.1137/18M1168285}.}
The generic odd-order question is distinct: the repository's TR-13 manuscript
uses a compressed Koszul flattening and explicitly leaves exceptional tensors
outside its conclusion.\footnote{M.~J.~Colbrook, repository manuscript for TR-13,
September 11, 2026. \url{https://github.com/ajt60gaibb/OpenProblemsInNLA/tree/main/tensor-computations/TR-13}.}
The exact lower bound below does not follow by specialization from a generic
rank statement.

Define the homogeneous moment functional
\[
 L_D(X^{D-j}Y^j)=h_j,\qquad 0\le j\le D.
\]
Equivalently, in the affine chart $X=1$, set $L(t^j)=h_j$ for $j\le D$.
Identifying a mode with $\Pq$, the tensor is the multilinear form
\begin{equation}\label{eq:moment}
 H(p_1,\ldots,p_m)=L(p_1\cdots p_m).
\end{equation}
Using dual spaces for this interpretation does not change either rank.
Vandermonde rank $\VR$ additionally requires each symmetric factor to lie on
$\{(1,t,\ldots,t^q):t\in\C\}\cup\{(0,\ldots,0,1)\}$, up to scaling.
Border rank, with any of these three restrictions, is the least rank bound
whose complex algebraic closure contains the tensor; equivalently, it permits
limits of tensors with that rank bound.
The associated binary form is
\[
 F(x,y)=\sum_{j=0}^D\binom Dj h_j x^{D-j}y^j.
\]
For $0\le d\le D$, define the apolar kernel
\begin{equation}\label{eq:apolarkernel}
 I_d=\left\{\sum_{i=0}^d g_iX^{d-i}Y^i:
              \sum_{i=0}^d g_i h_{i+j}=0\quad(0\le j\le D-d)\right\}.
\end{equation}
This is the usual degree-$d$ apolar kernel of $F$.
For $d>D$, put $I_d=\C[X,Y]_d$.
The spaces $I_d$ form a homogeneous ideal, directly from the displayed equations.

For $H\ne0$, let $r$ be the smallest degree for which $I_r\ne0$ and choose
$0\ne G\in I_r$. Let $s$ be the number of distinct zeros of $G$ in $\mathbb P^1$,
counted without multiplicity. Section~\ref{sec:momentalgebra} proves
\begin{equation}\label{eq:middle}
 r=\rank C_h,\qquad
 C_h=(h_{i+j})_{\substack{0\le i\le\lfloor D/2\rfloor\\
                         0\le j\le\lceil D/2\rceil}}.
\end{equation}
Thus the invariants in the next theorem are computable by exact linear algebra
and a univariate greatest common divisor.

\begin{theorem}[Exact rank formula; resolution of TR-14]\label{thm:main}
For every nonzero complex Hankel tensor of order $m\ge3$ and dimension $n\ge2$,
\begin{equation}\label{eq:main}
 \boxed{\displaystyle
 \Rk(H)=\SR(H)=
 \min\bigl\{D-r+2,\ (m-1)r-(m-2)s\bigr\}.}
\end{equation}
If $I_r$ has dimension greater than one, then $D=2r-2$ and both ranks are $r$,
regardless of the choice of $G$. The zero tensor has both ranks zero.
\end{theorem}

The proof has three ingredients. The minimal apolar polynomial gives a
nondegenerate finite moment algebra. Two symmetric constructions give the
upper bounds in \eqref{eq:main}. A product-space argument gives a matching
lower bound for an arbitrary minimum-length \emph{ordinary} decomposition.
That last argument, not a restriction of a full-algebra rank bound, is the
part that treats $r>n$ and multiple repeated support points.

\section{The finite moment algebra}\label{sec:momentalgebra}

A linear change in $(X,Y)$ acts invertibly on each degree-$q$ mode in
\eqref{eq:moment}. It therefore preserves ordinary and symmetric tensor rank,
and preserves the projective multiplicities of $G$. Choose such a change so
that $G$ has no zero at infinity. After scaling,
\[
 g(t)=G(1,t)=t^r+g_{r-1}t^{r-1}+\cdots+g_0
\]
is monic of degree $r$.
Write
\begin{equation}\label{eq:gfact}
 g(t)=\prod_{a=1}^s(t-\alpha_a)^{\ell_a},\qquad
 \ell_a\ge1,\quad \sum_a\ell_a=r.
\end{equation}

\begin{lemma}[Moment algebra]\label{lem:moment}
The algebra $A=\C[t]/(g)$ has a functional $\lambda:A\to\C$ such that
\begin{equation}\label{eq:lambda}
 \lambda([t^j])=h_j\quad(0\le j\le D).
\end{equation}
The pairing $(a,b)\mapsto\lambda(ab)$ is nondegenerate. Moreover,
$r\le\lfloor D/2\rfloor+1$, and \eqref{eq:middle} holds.
\end{lemma}
\begin{proof}
At degree $d=\lfloor D/2\rfloor+1$, the linear map in
\eqref{eq:apolarkernel} has $d+1$ columns and $D-d+1$ rows, with more columns
than rows. This proves the bound on $r$.
Define $\lambda$ on the basis $1,t,\ldots,t^{r-1}$ by the first $r$ moments.
The recurrence
\[
 h_{j+r}+g_{r-1}h_{j+r-1}+\cdots+g_0h_j=0
 \quad(0\le j\le D-r)
\]
then proves \eqref{eq:lambda} successively through degree $D$.

The radical of the pairing,
$J=\{a\in A:\lambda(aA)=0\}$, is an ideal. If $J\ne0$, then
$A/J\simeq\C[t]/(g')$ for a monic proper divisor $g'$ of $g$.
The functional descends to this quotient. Equation~\eqref{eq:lambda} implies
that $g'$ gives a recurrence for all the required moments, so its
homogenization is a nonzero element of $I_{\deg g'}$.
This contradicts the minimality of $r$. Hence the pairing is nondegenerate.

Both polynomial spaces of degrees $\lfloor D/2\rfloor$ and
$\lceil D/2\rceil$ surject onto $A$, because both degrees are at least $r-1$.
The middle catalecticant is the pullback of this nondegenerate pairing through
these two surjections. Its rank is therefore $r$. The rank is unchanged on
returning to the original binary coordinates.
\end{proof}

A functional with the nondegeneracy property in Lemma~\ref{lem:moment} is called
\emph{Frobenius}. The Chinese remainder decomposition is
\begin{equation}\label{eq:crt}
 A\simeq\prod_{a=1}^s A_a,
 \qquad A_a=\C[z_a]/(z_a^{\ell_a}),\quad z_a=t-\alpha_a.
\end{equation}
Every restriction $\lambda_a$ is Frobenius.
For a nonempty subset $J\subseteq\{1,\ldots,s\}$, let
$A_J=\prod_{a\in J}A_a$ and $r_J=\sum_{a\in J}\ell_a$.
Polynomial reduction gives
\begin{equation}\label{eq:restrictiondimensions}
 \dim \im(\C[t]_{\le q}\longrightarrow A_J)=\min\{n,r_J\}.
\end{equation}
Also,
\begin{equation}\label{eq:allbutone}
 \spanC\{[p_1]\cdots[p_{m-1}]:\deg p_i\le q\}=A,
\end{equation}
because $(m-1)q\ge r-1$ and products span every polynomial of degree at most
$(m-1)q$.

\begin{lemma}[Choice of the minimal apolar polynomial]\label{lem:unique}
If $D\ge2r-1$, then $I_r$ is one-dimensional. If $D=2r-2$, then
$\dim I_r=2$.
\end{lemma}
\begin{proof}
In the first case, polynomials of degree at most $D-r$ surject onto $A$.
For $b\in I_r$, the equations $\lambda([b][t^j])=0$ for $j\le D-r$ and
nondegeneracy give $[b]=0$. Since $\deg b\le r$, $b$ is a multiple of $g$.
In the second case, polynomials of degree at most $r-2$ give $r-1$ independent
functionals on $A$. Reduction of polynomials of degree at most $r$ onto $A$
has one-dimensional kernel, so the apolar kernel has dimension two.
\end{proof}

\section{Two symmetric upper bounds}\label{sec:upper}

\begin{proposition}[Local interpolation upper bound]\label{prop:localupper}
With \eqref{eq:gfact},
\begin{equation}\label{eq:localupper}
 \SR(H)\le\sum_{a=1}^s\bigl((m-1)(\ell_a-1)+1\bigr)
          =(m-1)r-(m-2)s.
\end{equation}
\end{proposition}
\begin{proof}
In $A_a=\C[z]/(z^\ell)$, write
\[
 \lambda_a(f)=[z^{\ell-1}]\,u(z)f(z),\qquad u(0)\ne0.
\]
Such $u$ exists uniquely, and its constant coefficient is nonzero because
$\lambda_a$ is Frobenius. Choose an $m$th root $w$ of $u$ modulo $z^\ell$;
this exists by the finite binomial expansion, after choosing an $m$th root
of $u(0)$.
Set $N=(m-1)(\ell-1)+1$. For polynomials $f_i$ of degree at most $\ell-1$,
\begin{equation}\label{eq:fourier}
 [z^{\ell-1}]\prod_{i=1}^m f_i(z)
 =\frac1N\sum_{\zeta^N=1}\zeta^{-(\ell-1)}
                    \prod_{i=1}^m f_i(\zeta).
\end{equation}
Indeed, among integers between $0$ and $m(\ell-1)$, the only one congruent
to $\ell-1$ modulo $N$ is $\ell-1$ itself.
Apply \eqref{eq:fourier} to
$f_i=w(z)p_i(\alpha_a+z)\bmod z^\ell$.
Each summand uses the same linear functional of $p_i$ in every mode, hence is
symmetric rank one. Summing the local decompositions gives \eqref{eq:localupper}.
The case $\ell=1$ has $N=1$ and is included.
\end{proof}

\begin{proposition}[Binary interpolation upper bound]\label{prop:binaryupper}
Every nonzero tensor under consideration satisfies
\begin{equation}\label{eq:binaryupper}
 \SR(H)\le\VR(H)\le D-r+2.
\end{equation}
\end{proposition}
\begin{proof}
Put $v=D-r+2\ge r$. Let $\tau:A\to\C$ take the coefficient of $t^{r-1}$
in the remainder modulo $g$. This is Frobenius: for a nonzero polynomial of
degree $d<r$, multiplication by $t^{r-1-d}$ produces a nonzero top coefficient.
Thus $\lambda(a)=\tau(ua)$ for a unique unit $u\in A$.
Choose the degree-at-most-$r-1$ representative $b(t)$ of $u^{-1}$.
For $0\le j\le r-2$,
\[
 L(t^j b(t))=\lambda([t^j b])=\tau([t^j])=0.
\]
Consequently the degree-$v$ homogenization $B(X,Y)=X^v b(Y/X)$ is apolar.
It is coprime to $G$: its residue modulo $g$ is a unit, and $G$ has no zero
at infinity.

Choose a homogeneous $Q$ of degree $v-r$ so that $GQ$ and $B$ are coprime
and $Q$ does not vanish at infinity. The pencil
$B+cGQ$ has a squarefree member with no root at infinity. To see this directly,
a repeated finite zero would be a critical point of the nonconstant rational
function $B/(GQ)$; in characteristic zero there are only finitely many
critical values. Exclude these finitely many values and $c=0$.

Let the resulting affine polynomial have distinct roots
$\beta_1,\ldots,\beta_v$. Solve the invertible Vandermonde system for weights
$\gamma_a$ matching $h_0,\ldots,h_{v-1}$. Apolarity gives the recurrence
which then matches every remaining moment through degree $D$. Hence
\[
 H=\sum_{a=1}^v\gamma_a(1,\beta_a,\ldots,\beta_a^q)^{\otimes m}.
\]
Zero weights may be discarded. For $r=1$, the same argument uses an empty
set of conditions $0\le j\le-1$; alternatively Lemma~\ref{lem:moment}
immediately gives rank one.
\end{proof}

\begin{proposition}[Vandermonde rank dichotomy]\label{prop:vandermonde}
The exact Vandermonde rank is
\[
 \VR(H)=\begin{cases}
 r,&G\text{ is squarefree},\\
 D-r+2,&G\text{ is not squarefree}.
 \end{cases}
\]
This is independent of the choice of $G$ in the balanced case $D=2r-2$.
\end{proposition}
\begin{proof}
A Vandermonde decomposition with $d$ distinct projective nodes gives a
squarefree polynomial in $I_d$: take the product of the linear forms vanishing
at the nodes. Conversely, a squarefree apolar polynomial yields a decomposition
on its nodes by the Vandermonde recurrence construction, after moving infinity
away from the nodes. Thus $\VR(H)$ is the smallest squarefree apolar degree.

For $r\le d<D-r+2$, every $b\in I_d$ satisfies
$\lambda([b][t^j])=0$ for $j\le D-d$, and $D-d\ge r-1$.
Nondegeneracy forces $[b]=0$, so in homogeneous notation
$I_d=G\,\C[X,Y]_{d-r}$. If $G$ is not squarefree, none of these polynomials
is squarefree; Proposition~\ref{prop:binaryupper} attains degree $D-r+2$.
If $G$ is squarefree, degree $r$ is already attainable and minimal by definition
of $r$. When $D=2r-2$, both alternatives give $r$.
\end{proof}

The familiar binary interpolation construction is thus sufficient for the
first bound; no genericity is used. The local construction supplies the
potentially smaller second bound even when $r>n$.

\section{A product-space inequality with a Frobenius constraint}\label{sec:product}

For subspaces $U,V$ of a commutative algebra, $UV$ denotes the linear span of
all products $uv$. Empty products of subspaces mean $\C1$.
For a subspace $W$, define
\[
 \Stab(W)=\{b:bW\subseteq W\}.
\]
It is a unital subalgebra.
The following transformation lemma is the commutative finite-dimensional
specialization of Beck--Lecouvey, Lemma 4.2.\footnote{V.~Beck and C.~Lecouvey,
\emph{Additive combinatorics methods in associative algebras}, Confluentes
Math. 9 (2017), no. 1, 3--27, Lemma 4.2.
\url{https://doi.org/10.5802/cml.34}.}

\begin{lemma}[Linear transformation lemma]\label{lem:transform}
Let $B$ be a finite-dimensional commutative complex algebra. If $U,V\subseteq B$
contain units, then for every unit $a\in U$ there are a unital subalgebra
$K_a\subseteq B$ and a subspace $T_a$ such that
\begin{equation}\label{eq:transform}
 aV\subseteq T_a\subseteq UV,\qquad K_aT_a=T_a,\qquad
 \dim T_a+\dim K_a\ge\dim U+\dim V.
\end{equation}
\end{lemma}
\begin{proof}
Scaling reduces to $a=1\in U$ and $1\in V$. Induct on $\dim U$.
For a unit $e\in V$, put
\[
 U_e=U\cap Ve^{-1},\qquad V_e=V+Ue.
\]
Then $1\in U_e\cap V_e$, $U_eV_e\subseteq UV$, and
$\dim U_e+\dim V_e=\dim U+\dim V$.
If some $U_e\ne U$, apply induction to $(U_e,V_e)$.
Otherwise $Ue\subseteq V$ for every unit $e\in V$.
Units span $V$ (the nonunits form finitely many proper residue hyperplanes),
so $UV=V$. The algebra $K=\C[U]$ satisfies $KV=V$ and $\dim K\ge\dim U$;
take $T=V$. This also covers the initial case $\dim U=1$.
Undoing the unit scalings gives \eqref{eq:transform}.
\end{proof}

A general product-space Kneser inequality cannot simply be invoked in an
arbitrary nonreduced algebra. The next lemma proves exactly the contextual
version needed here and supplies the missing finiteness argument.

\begin{lemma}[Contextual product inequality]\label{lem:context}
Let $B=A\times\C^R$, where $A$ is a finite-dimensional commutative complex
algebra. Let $\Lambda:B\to\C$ restrict to a Frobenius functional on $A$.
Suppose $U_1,\ldots,U_m\subseteq B$ contain $1$, and set $W=U_1\cdots U_m$.
Assume
\begin{equation}\label{eq:contextassumptions}
 \Lambda(W)=0,\qquad
 \pi_A\left(\prod_{k\ne i}U_k\right)=A
 \quad\text{for every }i.
\end{equation}
For $W_j=U_1\cdots U_j$ and $j\ge2$,
\begin{equation}\label{eq:knesercontext}
 \dim W_j\ge\dim W_{j-1}+\dim U_j-\dim\Stab(W_j).
\end{equation}
Moreover, $\Stab(W)$ is reduced.
\end{lemma}
\begin{proof}
Fix $j$, put $U=U_j$, $V=W_{j-1}$, and
$Q=U_{j+1}\cdots U_m$.
Apply Lemma~\ref{lem:transform} to a unit $a\in U$, obtaining $K_a,T_a$.
If $z\in K_a$ is nilpotent, its $\C^R$ component is zero. Since
$zT_a\subseteq T_a$,
\[
 z\,aVQ\subseteq T_aQ\subseteq UVQ=W.
\]
Consequently
\[
 \Lambda_A(z_Aa_A\,\pi_A(VQ))=0.
\]
Here $\pi_A(VQ)=A$ by \eqref{eq:contextassumptions}, and $a_A$ is a unit.
Nondegeneracy of the Frobenius pairing forces $z_A=0$, hence $z=0$.
Thus \emph{every auxiliary algebra $K_a$ is reduced}.

There are only finitely many reduced unital subalgebras of $B$.
Indeed, a reduced finite-dimensional complex algebra is a product of copies
of $\C$. Its primitive idempotents map to idempotents of $B$.
In each local factor of $B$ an idempotent is either $0$ or $1$.
Such subalgebras are therefore specified by partitions of the finite set of
local factors of $B$.

Let $d=\dim U$ and choose a basis $x_1=1,x_2,\ldots,x_d$ of $U$.
For $c\in\C$, set
\[
 a(c)=x_1+cx_2+\cdots+c^{d-1}x_d.
\]
All but finitely many $a(c)$ are units: each residue coordinate is a polynomial
with constant term $1$. Finiteness of the possible $K_{a(c)}$ gives $d$ distinct
values $c_1,\ldots,c_d$ with a common auxiliary algebra $K$.
The corresponding $a(c_i)$ form a basis of $U$ by the Vandermonde determinant.
It follows that
\[
 UV=\sum_{i=1}^d a(c_i)V=\sum_{i=1}^dT_{a(c_i)},\qquad KUV=UV.
\]
The first equality with the $T$-spaces follows by the two inclusions in
\eqref{eq:transform}. Therefore $K\subseteq\Stab(UV)$, and
\[
 \dim UV+\dim\Stab(UV)
 \ge\dim T_{a(c_1)}+\dim K\ge\dim U+\dim V.
\]
This proves \eqref{eq:knesercontext}.

Finally, if $z\in\Stab(W)$ is nilpotent, then $z_{\C^R}=0$ and
$\Lambda_A(z_A\pi_A(W))=0$. The projection $\pi_A(W)=A$, since every $U_i$
contains $1$ and \eqref{eq:contextassumptions} holds.
Again nondegeneracy gives $z=0$, proving that $\Stab(W)$ is reduced.
\end{proof}

\begin{corollary}[Trivial-stabilizer bound]\label{cor:trivial}
Under the hypotheses of Lemma~\ref{lem:context}, if $\Stab(W)=\C1$, then
\begin{equation}\label{eq:trivial}
 \dim W\ge\sum_{i=1}^m\dim U_i-(m-1).
\end{equation}
\end{corollary}
\begin{proof}
Every $\Stab(W_j)$ is contained in $\Stab(W)$, because multiplication by the
remaining factors preserves its stabilizing action. Its dimension is therefore
one. Sum \eqref{eq:knesercontext} for $j=2,\ldots,m$.
\end{proof}

\section{Lower bound for an arbitrary ordinary decomposition}\label{sec:lower}

\begin{proposition}[Exact ordinary-rank lower bound]\label{prop:lower}
For the moment tensor in \eqref{eq:moment},
\begin{equation}\label{eq:lower}
 \Rk(H)\ge\min\{D-r+2,\ (m-1)r-(m-2)s\}.
\end{equation}
\end{proposition}
\begin{proof}
Let $R=\Rk(H)$ and take a minimum-length ordinary decomposition,
\begin{equation}\label{eq:ordinary}
 H(p_1,\ldots,p_m)=\sum_{b=1}^R\prod_{i=1}^m\beta_{b,i}(p_i),
 \qquad \beta_{b,i}\in(\Pq)^*\setminus\{0\}.
\end{equation}
Its $R$ decomposable tensors are linearly independent. Otherwise a nonzero
linear relation would eliminate one summand after rescaling the others,
contradicting minimality.

Set $B=A\times\C^R$ and introduce the graph subspaces
\[
 \widetilde U_i=
 \{([p],\beta_{1,i}(p),\ldots,\beta_{R,i}(p)):p\in\Pq\}\subseteq B.
\]
Define
\[
 \widetilde\Lambda(a,c_1,\ldots,c_R)=\lambda(a)-\sum_{b=1}^R c_b.
\]
Equation~\eqref{eq:ordinary} says
$\widetilde\Lambda(\widetilde U_1\cdots\widetilde U_m)=0$.
Choose a unit $u_i\in\widetilde U_i$ for every $i$.
This is possible by avoiding the finitely many hyperplanes
$p(\alpha_a)=0$ and $\beta_{b,i}(p)=0$.
Now put
\begin{equation}\label{eq:normalize}
 U_i=u_i^{-1}\widetilde U_i,\qquad
 \Lambda(x)=\widetilde\Lambda(u_1\cdots u_mx),\qquad
 W=U_1\cdots U_m.
\end{equation}
Then $1\in U_i$, $\Lambda(W)=0$, the restriction $\Lambda_A$ is Frobenius,
and each coefficient of $\Lambda$ on a $\C$ factor is nonzero.
Equation~\eqref{eq:allbutone}, multiplied by units, verifies every surjectivity
hypothesis in \eqref{eq:contextassumptions}.

In addition,
\begin{equation}\label{eq:ontoB}
 \pi_{\C^R}(W)=\C^R.
\end{equation}
Indeed, the coordinate functionals of the unnormalized product-evaluation map
are the linearly independent decomposable tensors in \eqref{eq:ordinary}.
The map is therefore onto; normalization is just invertible coordinate scaling.

By Lemma~\ref{lem:context}, $K=\Stab(W)$ is reduced. Write its primitive
idempotents as $e_1,\ldots,e_t$, so
\[
 B=\prod_{j=1}^t B_j,\qquad B_j=e_jB=A_j\times\C^{R_j},\qquad
 W=\bigoplus_{j=1}^t W^{(j)},\quad W^{(j)}=e_jW.
\]
Each $A_j$ is a product of some of the local factors in \eqref{eq:crt}.
Write $r_j=\dim A_j$ and let $s_j$ count those local factors.
No $A_j$ is zero. Otherwise \eqref{eq:ontoB} would give
$W^{(j)}=B_j=\C^{R_j}$, contradicting $\Lambda(W^{(j)})=0$ and the nonzero
coordinate coefficients of $\Lambda$. Hence
\begin{equation}\label{eq:partitions}
 r_j\ge1,\quad s_j\ge1,\quad t\le s,\quad
 \sum_j r_j=r,\quad\sum_jR_j=R.
\end{equation}

Put $U_i^{(j)}=e_jU_i$. These spaces contain the identity $e_j$ of $B_j$,
and their product is $W^{(j)}$. The contextual hypotheses pass to $B_j$.
The stabilizer of $W^{(j)}$ inside $B_j$ is exactly $\C e_j$:
an additional stabilizer, extended by zero to the other blocks, would be an
element of $K$ not constant on its $j$th primitive block.
By \eqref{eq:restrictiondimensions} and normalization by units,
\[
 \dim U_i^{(j)}\ge\min\{n,r_j\}.
\]
Corollary~\ref{cor:trivial} now yields
\[
 \dim W^{(j)}\ge m\min\{n,r_j\}-m+1.
\]
On the other hand, the nonzero functional $\Lambda|_{B_j}$ annihilates
$W^{(j)}$, so $\dim W^{(j)}\le r_j+R_j-1$. Therefore
\begin{equation}\label{eq:blockbound}
 \boxed{\displaystyle R_j\ge m\min\{n,r_j\}-r_j-m+2.}
\end{equation}

There are two possibilities. If some $r_j\ge n$, then
\[
 R\ge R_j\ge mn-r_j-m+2=D-r_j+2\ge D-r+2.
\]
If every $r_j<n$, sum \eqref{eq:blockbound} and use \eqref{eq:partitions}:
\[
 R\ge\sum_j((m-1)r_j-(m-2))
   =(m-1)r-(m-2)t
   \ge(m-1)r-(m-2)s.
\]
In either case \eqref{eq:lower} follows.
\end{proof}

\begin{proof}[Proof of Theorem~\ref{thm:main}]
Propositions~\ref{prop:localupper} and \ref{prop:binaryupper} give the upper
bound for $\SR(H)$, Proposition~\ref{prop:lower} gives the identical lower
bound for $\Rk(H)$, and always $\Rk(H)\le\SR(H)$.
If $I_r$ is not one-dimensional, Lemma~\ref{lem:unique} gives $D=2r-2$.
Then $D-r+2=r$ while $(m-1)r-(m-2)s\ge r$, so the result is independent of $G$.
\end{proof}

\begin{remark}[Why the mixed-root case is included]
The idempotents in the argument partition whole local factors of $A$, not
individual derivative coordinates. Thus a repeated root cannot be silently
split into reduced points. The graph subspaces record arbitrary mode-dependent
factors $\beta_{b,i}$; none is assumed symmetric or Vandermonde.
The algebra $A$ can be larger than a mode space. Only products of all but one
mode are required to surject onto $A$, as verified in \eqref{eq:allbutone}.
\end{remark}

\section{Border ranks and a fixed Koszul identity}\label{sec:border}

\begin{theorem}[All border ranks]\label{thm:border}
For every complex Hankel tensor,
\begin{equation}\label{eq:border}
 \BR(H)=\BSR(H)=\BVR(H)=\rank C_h=r.
\end{equation}
\end{theorem}
\begin{proof}
The zero tensor is immediate. For a nonzero tensor, work in the monic
coordinates of Section~\ref{sec:momentalgebra} for the upper bound.
Perturb $g$ to squarefree monic polynomials $g_\varepsilon$ of degree $r$.
Keep the first $r$ moments fixed and extend them through degree $D$ by the
$g_\varepsilon$ recurrence. The resulting moments converge to the original
ones, and the Vandermonde system gives a length-at-most-$r$ Vandermonde
decomposition for each perturbation. This proves $\BVR(H)\le r$.
It remains to prove the unrestricted border-rank lower bound.

For even $m=2k$, group the modes into two groups of $k$. For each possible
index sum choose one representative multi-index in each group.
The corresponding submatrix of the flattening is $C_h$, so its rank gives
$\BR(H)\ge r$.

Now suppose $m=2k+1$. Put
\[
 a=kq+1,\quad b=\lfloor q/2\rfloor,\quad c=\lceil q/2\rceil,
 \quad p_1=a+b,\quad p_2=a+c.
\]
Group the modes as $k,1,k$. In the two outer groups, extract one coordinate
for each sum $0,\ldots,kq$. In the middle mode, extract coordinates $0,b,q$.
These are linear maps, even when $q=1$ and two extracted coordinates coincide.
The resulting $a\times3\times a$ tensor has slices
\[
 M_j=(h_{u+v+j})_{0\le u,v<a},\qquad j=0,b,q.
\]
Its three-slice Koszul matrix is
\begin{equation}\label{eq:K}
 \mathcal K=\begin{pmatrix}
 0&M_q&-M_b\\-M_q&0&M_0\\M_b&-M_0&0
 \end{pmatrix}.
\end{equation}
On an arbitrary rank-one $a\times3\times a$ tensor this matrix has rank at
most two: it is the tensor product of a $3\times3$ skew matrix of rank two
and a rank-one matrix. Therefore
$\rank\mathcal K\le2\BR(H)$, by linearity and closedness of matrix-rank bounds.

Here $C_h$ is $p_1\times p_2$. For $x\in\C^{p_1}$ and $y\in\C^{p_2}$,
define the $(3a)\times(p_1+p_2)$ matrix $E$ by
\begin{equation}\label{eq:E}
 E(x,y)=\bigl((x_{i+b})_{i=0}^{a-1},\
                  (-x_i+y_{i+c})_{i=0}^{a-1},\
                  (-y_i)_{i=0}^{a-1}\bigr).
\end{equation}
A direct block multiplication, using $b+c=q$, gives the identity
\begin{equation}\label{eq:identity}
 \boxed{\displaystyle
 \mathcal K=E\begin{pmatrix}0&C_h\\-C_h^{\mathsf T}&0\end{pmatrix}E^{\mathsf T}.}
\end{equation}
It holds for every moment vector, not only a generic one.

The matrix $E$ has full column rank. If $E(x,y)=0$, then
$x_i=0$ for $i\ge b$, and $y_i=0$ for $i<a$.
For $i<b$, the middle equations give $x_i=y_{i+c}=0$, since $i+c<q\le a$.
Thus $x=0$. All remaining $y$ coordinates are then zero by the same middle
equations, because $a\ge c$.
Consequently \eqref{eq:identity} gives
$\rank\mathcal K=2\rank C_h=2r$.
This proves $\BR(H)\ge r$ also for odd order, and proves \eqref{eq:border}.
\end{proof}

\section{Consequences and examples}\label{sec:examples}

\begin{corollary}[Reduced support]
If $G$ is squarefree, then all ordinary, symmetric, and Vandermonde exact and
border ranks equal $r$.
\end{corollary}
\begin{proof}
Here $s=r$, so Theorem~\ref{thm:main} gives exact rank $r$.
The same distinct-root recurrence gives a Vandermonde decomposition of length
$r$, and Theorem~\ref{thm:border} gives the lower bound.
\end{proof}

\begin{corollary}[A single support point]
If $G$ has just one projective zero, then
\[
 \Rk(H)=\SR(H)=\min\{D-r+2,\ (m-1)(r-1)+1\}.
\]
In particular, if $h_k\ne0$ and $h_j=0$ for all $j>k$, with $k\le D/2$, then
$r=k+1$, $s=1$, and
\[
 \Rk(H)=\SR(H)=\min\{D-k+1,\ (m-1)k+1\}.
\]
\end{corollary}
\begin{proof}
The first assertion substitutes $s=1$. In the second, $t^{k+1}$ gives a
recurrence and the $(k+1)\times(k+1)$ leading Hankel block is anti-triangular
with nonzero anti-diagonal $h_k$. Hence its rank is $k+1$, giving the stated
minimal polynomial and rank.
\end{proof}

\begin{corollary}[Maximum Hankel rank]\label{cor:max}
For fixed $m,n$,
\[
 \max_{H\ \mathrm{Hankel}}\Rk(H)
 =\max_{H\ \mathrm{Hankel}}\SR(H)=(m-1)(n-1)+1.
\]
\end{corollary}
\begin{proof}
If $r\le n$, the second term of \eqref{eq:main} is at most
$(m-1)n-(m-2)=(m-1)q+1$, since $s\ge1$.
If $r\ge n$, the first term is at most $D-n+2=(m-1)q+1$.
The moment sequence $h_q=1$, $h_j=0$ for $j\ne q$, has $r=n$, $s=1$ and
attains the bound.
\end{proof}

\begin{example}[A mixed repeated-root case with $r>n$]\label{ex:mixed}
Take $m=5$, $n=3$, and
\[
 L(f)=f'(0)+f(1)+f(2),\qquad
 h_j=\mathbf1_{\{j=1\}}+1+2^j\quad(0\le j\le10).
\]
The minimal polynomial is $g(t)=t^2(t-1)(t-2)$, so $r=4$, $s=3$.
The formula gives
\[
 \Rk(H)=\SR(H)=\min\{8,7\}=7,
 \qquad \BR(H)=4.
\]
The upper bound is visible in the exact decomposition
\[
 H=\frac15\sum_{\zeta^5=1}\zeta^{-1}(1,\zeta,0)^{\otimes5}
       +(1,1,1)^{\otimes5}+(1,2,4)^{\otimes5}.
\]
The lower bound is Proposition~\ref{prop:lower}; it does not assume that an
ordinary decomposition has the displayed form. This example distinguishes
exact rank, border rank, and the degree-eight binary interpolation bound.
\end{example}

\begin{example}[Two repeated support points above the mode dimension]
Take $m=5$, $n=3$, and $L(f)=f'(0)+f'(1)$.
Then $g=t^2(t-1)^2$, $r=4$, $s=2$. The local construction gives ten terms,
whereas the binary construction gives eight. Theorem~\ref{thm:main} proves
that both exact ranks are eight.
\end{example}

\begin{example}[Rank substantially below Vandermonde rank]
For $m\ge3$, $n=3$, let $h_1=h_{2m}=1$ and all other moments vanish.
The associated form in three tensor variables is
$m x^{m-1}y+z^m$.
Its projective apolar support has one double point and one simple point:
$r=3$, $s=2$. Thus
\[
 \Rk(H)=\SR(H)=m+1,
 \qquad \BR(H)=3.
\]
Its binary Vandermonde rank is $2m-1$, by Proposition~\ref{prop:vandermonde}. Equality of ordinary and symmetric rank therefore does not mean
that an optimal decomposition must stay on the rational normal curve.
\end{example}

\begin{example}[A table of exact values]
Each row specifies a Frobenius moment functional by the multiplicities of its
minimal apolar support. The balanced row is independent of that choice.
\begin{center}
\begin{tabular}{@{}rrrrrrr@{}}
\toprule
$m$&$n$&$r$&$s$&$D-r+2$&$(m-1)r-(m-2)s$&$\Rk=\SR$\\
\midrule
3&3&3&1&5&5&5\\
3&4&5&4&6&6&6\\
4&3&4&2&6&8&6\\
5&3&4&3&8&7&7\\
5&3&4&2&8&10&8\\
5&3&6&1&6&21&6\\
7&4&5&4&18&10&10\\
\bottomrule
\end{tabular}
\end{center}
\end{example}

\section{An exact rank algorithm}\label{sec:algorithm}

For exact rational or algebraic input, no root approximation is required.
Form the matrices
\[
 C_d=(h_{i+j})_{\substack{0\le j\le D-d\\0\le i\le d}}
 \qquad(d=1,\ldots,\lfloor D/2\rfloor+1).
\]
The first matrix with nonzero kernel gives $r$ and coefficients
$(g_0,\ldots,g_r)$ of $G=\sum g_iX^{r-i}Y^i$.
Let $g(t)=\sum g_it^i$, and let $e=\deg g$ in the original coordinates.
The number of distinct projective zeros is
\begin{equation}\label{eq:gcdcount}
 s=e-\deg\gcd(g,g')+\mathbf1_{\{e<r\}}.
\end{equation}
For a nonzero constant $g$, the first two terms are zero.
The final indicator counts the possible root at infinity only once, irrespective
of its multiplicity. Return the minimum in \eqref{eq:main}.

The supplied \texttt{code/hankel\_rank.py} implements these operations in
SymPy using exact rational arithmetic. It reports the two upper bounds, the
apolar polynomial, the ordinary and symmetric ranks, and the common border
rank. It rejects floating-point input rather than deciding exact rank with a
numerical tolerance. Rank computation does not require factoring $G$ into
linear factors.

The companion \texttt{code/binary\_certificate.py} constructs a minimal
Vandermonde decomposition implicitly, without approximating its algebraic roots.
Its output consists of a squarefree monic polynomial $P$, a polynomial $U$,
and a vector of polynomials $f$, all over $\mathbb Q$, with
\[
 H=\sum_{P(\beta)=0}\frac{U(\beta)}{P'(\beta)}f(\beta)^{\otimes m}.
\]
The top-coefficient identity for Lagrange interpolation verifies this sum by
polynomial remainders in $\mathbb Q[t]/(P)$. A rational projective shear handles
roots initially at infinity. These Vandermonde certificates are optimal for
ordinary and symmetric rank precisely when their length agrees with
\eqref{eq:main}; the output explicitly labels this distinction.

An optimal symmetric decomposition can also be constructed. If the second
term of \eqref{eq:main} is smaller, use the Chinese remainder and root-of-unity
construction in Proposition~\ref{prop:localupper}. If the first term is smaller,
use the squarefree apolar pencil in Proposition~\ref{prop:binaryupper}.
Any zero or redundant summands would contradict the lower bound, so a
length-achieving construction is minimal.

\section{Verification scope}\label{sec:verification}

The mathematical proof of the universal statement is Sections~2--5.
The finite computations supplied with this manuscript are reproducibility
checks, not a replacement for that proof. They check exact catalecticant and
apolar calculations, local Fourier interpolation identities, complete
small-dimensional tensor reconstructions, a mixed-root graph-algebra example,
and the fixed identity \eqref{eq:identity} across many dimensions and moment
patterns. Finite-field checks are labeled as such and do not purport to prove
the complex-field theorem by testing.

A separate Codex AI agent independently audited this complete argument on
12 September 2026 and returned PASS under the repository's informal-audit
policy. The dated review is available at
\url{https://github.com/sidneyholden1/OpenProblemsInNLA/blob/codex/holden-tr14-resolution/references/holden-tr14-2026-09-12/independent-review.md}.
AI assistance was used for review and submission preparation. This does not
constitute external human peer review or formal verification; no Lean checking
was performed. In particular, a successful run of the verification scripts is not
a certificate for an arbitrary claimed mathematical proof. The companion
\texttt{PROOF\_AUDIT.md} identifies the dependencies and checks the
surjectivity, nilpotent exclusion, and support-partition steps explicitly.

\appendix
\section{An independent lower-bound check for terminal moment sequences}

This appendix provides an elementary verification of the single-support
formula without the contextual product-space argument.

\begin{lemma}[Terminal-sequence substitution bound]\label{lem:terminal}
Let $m\ge2$, $q\ge1$, $q\le k\le mq$, and let $T$ have entries
$h_{i_1+\cdots+i_m}$, $0\le i_j\le q$, where $h_k\ne0$ and $h_j=0$ for $j>k$.
Then $\Rk(T)\ge mq-k+1$.
\end{lemma}
\begin{proof}
Induct on $m$. For $m=2$, an anti-triangular submatrix of size $2q-k+1$
has nonzero anti-diagonal $h_k$.
For general $m$, if $k>(m-1)q$, the first-mode flattening contains an
anti-triangular minor of size $mq-k+1$: use row indices
$i=k-(m-1)q,\ldots,q$ and column sums $k-q,\ldots,(m-1)q$.

It remains that $q\le k\le(m-1)q$. The first-mode flattening has full row
rank $q+1$, using rows $0,\ldots,q$ and column sums $k-q,\ldots,k$.
In a minimum decomposition, choose $q+1$ independent first-mode factors and
their dual basis $p_0,\ldots,p_q$ in $\C[t]_{\le q}$.
At least one $p_a$ has $p_a(0)\ne0$. Contraction with that polynomial kills
$q$ of the chosen summands and gives an $(m-1)$-mode Hankel tensor with moments
$h'_j=\sum_{u=0}^q(p_a)_u h_{j+u}$.
It has $h'_k=(p_a)_0h_k\ne0$ and $h'_j=0$ for $j>k$.
Induction gives its rank at least $(m-1)q-k+1$; restoring the $q$ killed
summands gives $mq-k+1$.
\end{proof}

For $1\le k\le q$, all entries with a mode index above $k$ vanish. Restrict
each mode to indices $0,\ldots,k$ and apply Lemma~\ref{lem:terminal} with
$q'=k$; this gives $(m-1)k+1$. The case $k=0$ has rank one directly.
For $k\ge q$, the lemma supplies $D-k+1$. When $k\le D/2$, these lower bounds
match the two branches of the single-support formula and the symmetric upper
bounds in Section~3. Thus no contextual product inequality is needed for this
check, including the extremal example in Corollary~\ref{cor:max}.

\section{Source notes}

The contextual inequality in Lemma~\ref{lem:context}, and its use with the graph
of an unrestricted decomposition, are the essential arguments of this
manuscript. The linear transformation in Lemma~\ref{lem:transform} is credited
to Beck--Lecouvey. Its finiteness conclusion here is not assumed from a theorem
requiring finitely many subalgebras of the ambient nonreduced algebra; it is
proved by excluding nilpotents from each auxiliary algebra.

For $r\le n$, the lower bound reduces to the familiar commutative
Alder--Strassen-type expression. The general multilinear algebra lower bound
provides a useful comparison, but is not used to pass a rank bound through a
restriction.\footnote{M.~Bl\"aser, H.~Mayer, and D.~Shringi,
\emph{On the Multilinear Complexity of Associative Algebras}, STACS 2023,
Article 12, Theorem 19. \url{https://doi.org/10.4230/LIPIcs.STACS.2023.12}.}

\begin{thebibliography}{9}
\bibitem{TR14}
Open Problems in Numerical Linear Algebra.
\emph{TR-14: Comon's exact-rank conjecture for Hankel tensors}.
Repository page, last checked there September 10, 2026; accessed September 12, 2026.
\url{https://github.com/ajt60gaibb/OpenProblemsInNLA/tree/main/tensor-computations/TR-14}.

\bibitem{NY}
J.~Nie and K.~Ye.
\emph{Hankel Tensor Decompositions and Ranks}.
SIAM J. Matrix Anal. Appl. \textbf{40} (2019), 486--516.
\url{https://doi.org/10.1137/18M1168285}.
Author preprint: \url{https://arxiv.org/abs/1706.03631}.

\bibitem{BL}
V.~Beck and C.~Lecouvey.
\emph{Additive combinatorics methods in associative algebras}.
Confluentes Math. \textbf{9} (2017), no. 1, 3--27.
\url{https://doi.org/10.5802/cml.34}.
Author preprint: \url{https://arxiv.org/abs/1504.02287}.

\bibitem{TR13}
M.~J.~Colbrook.
\emph{TR-13: Equality of ranks for generic odd-order Hankel tensors}.
Repository manuscript, September 11, 2026.
\url{https://github.com/ajt60gaibb/OpenProblemsInNLA/tree/main/tensor-computations/TR-13}.
The repository describes this as recovered AI-assisted work, not external
human peer review or formal certification.

\bibitem{BMS}
M.~Bl\"aser, H.~Mayer, and D.~Shringi.
\emph{On the Multilinear Complexity of Associative Algebras}.
40th International Symposium on Theoretical Aspects of Computer Science,
LIPIcs \textbf{254} (2023), Article 12, 12:1--12:18.
\url{https://doi.org/10.4230/LIPIcs.STACS.2023.12}.
\end{thebibliography}
\end{document}
